\section{What remains}
\label{sec:outlook}

Theorem~\ref{thm:main} is the first unconditional step; this section names the next two, sharpening the power and widening the range of $\sigma$, and the inequality each of them rests on.

\paragraph*{The power.}
Theorem~\ref{thm:main} has the power $\ell^{-5/4}(\ln\ell)^{3/4}$; the conjecture~\eqref{eq:conj1} has $\ell^{-1/2}$.
The whole difference sits in Theorem~\ref{thm:circles}.
With $\keff^2\le c_0\ln(1/A)+C_0$ the same chain gives $\ell^{-1/2}$.
By Lemma~\ref{lem:smooth} and the half lemma, this sharp bound is equivalent to a stability of the $w$-information under a lowering of the $v$-noise:
\begin{equation}
F\big(1-\tfrac1{1+\ln(1/A)}\big)\le\rho\,F(1)\quad\text{for a universal }\rho .
\label{eq:S}
\end{equation}
On every recorded state the ratio $\rho_*=F(1-1/(1+L))/F(1)$ is below $e$: at most $2.700$ on $2921$ one-variable states and $2.648$ on $2449$ two-dimensional ones, the largest values coming from valley mixtures, tori and lattices \numst.
We conjecture that~\eqref{eq:S} holds with $\rho=e$ \conjst.
The three circles prove~\eqref{eq:S} only for the smaller step $\tau=\ln(2L)/(8L^2)$ and with the factor $36.1\sqrt{L/\ln2L}$.
The obstacle is the return from the weight $1/\varrho_T$ to $1/\varrho_{1-\tau}$ at a $v$-noise of order one (Remark~\ref{rem:delta}).
A version that separates the far, heavy parts of the state from its local ripple at each point is the natural next attempt.

\paragraph*{The range.}
Table~\ref{tab:range} shows that better constants widen the range only slowly.
Two changes of form look more promising.
(i) A relative bound on $\Wzz$ inside $\Oset$, in place of the absolute smallness $J\le J_0$.
(ii) Writing the region outside $\Oset$ also with the expansion of Lemma~\ref{lem:EGD}.
Its remainder carries $\sqrt J$ instead of the state-independent $(\sqrt{960}/6)\sigma^2$, so the threshold $112\sigma^2$ of Lemma~\ref{lem:EF2}, and with it the constant $1268$ of Lemma~\ref{lem:EFO}, could be lowered or removed.

\paragraph*{The upper side and the window class.}
This paper is about the lower side.
The upper side of the logarithmic law, and the coefficient one of Papers~I and~J for windows, are not touched.
